% Generated draft: requires mathematical and source review.
\documentclass[11pt]{article}
\usepackage[margin=25mm]{geometry}
\usepackage{fontspec}
\setmainfont{Latin Modern Roman}
\usepackage{amsmath,amssymb,mathtools,unicode-math}
\setmathfont{Latin Modern Math}
\usepackage{enumitem,needspace,xurl,hyperref}
\setlength{\parindent}{0pt}
\setlength{\parskip}{5pt}
\setlength{\emergencystretch}{3em}
\newcommand{\N}{\mathbb{N}}
\newcommand{\Z}{\mathbb{Z}}
\newcommand{\Q}{\mathbb{Q}}
\newcommand{\R}{\mathbb{R}}
\newcommand{\C}{\mathbb{C}}
\newcommand{\F}{\mathbb{F}}
\newcommand{\A}{\mathbb{A}}
\newcommand{\PP}{\mathbb{P}}
\newcommand{\E}{\mathbb{E}}
\newcommand{\eps}{\varepsilon}
\newcommand{\dd}{\,\mathrm{d}}
\begin{document}
\section*{UnsolvedMath ID: 30001542}
\textbf{Source problem code:} catalogue entry
\subsection{Definitions and mathematical statement}
The supplied question concerns surface types denoted by \(( -\chi,p)\), where \(\chi\) is the Euler characteristic. The record does not define \(p\), specify the class of surfaces, or state the conventions for hyperbolicity and trace fields. Consequently, the following is a conditional formulation under explicitly stated standard conventions; these conventions require confirmation against the cited report.

Suppose that \(S_{g,p}\) denotes a connected orientable surface obtained by deleting \(p\) points from a closed surface of genus \(g\), where \(g,p\in\Z_{\geq 0}\). Thus
\[
\chi(S_{g,p})=2-2g-p.
\]
For an orientation-preserving homeomorphism \(f:S_{g,p}\to S_{g,p}\), define its mapping torus by
\[
M_f=(S_{g,p}\times[0,1])/\bigl((x,1)\sim(f(x),0)\bigr).
\]
In this conditional formulation, a hyperbolic mapping torus means that \(M_f\) admits a complete hyperbolic metric of finite volume. For such a metric, let
\[
\rho:\pi_1(M_f)\longrightarrow\operatorname{PSL}_2(\C)
\]
be its holonomy representation, let \(q:\operatorname{SL}_2(\C)\to\operatorname{PSL}_2(\C)\) be the quotient map, and define the ordinary trace field by
\[
K(M_f)=\Q\bigl(\operatorname{tr}(A):A\in q^{-1}(\rho(\pi_1(M_f)))\bigr).
\]
A real place here means a field embedding \(K(M_f)\hookrightarrow\R\).

Under these conventions, the question is whether
\[
\begin{gathered}
\text{for every }g,p\in\Z_{\geq0}\text{ with }2g-2+p>1,\\
\text{there exists an orientation-preserving homeomorphism }f:S_{g,p}\to S_{g,p}\\
\text{such that }M_f\text{ admits a complete finite-volume hyperbolic metric}\\
\text{and }K(M_f)\text{ admits an embedding into }\R.
\end{gathered}
\]
The source quantifies over every surface type \(( -\chi,p)\) with \(-\chi>1\). Its identification with the genus-and-puncture formulation above is conditional on the unstated conventions. In particular, punctures versus boundary components, orientability, finite volume, and ordinary versus invariant trace field must be checked. The background also refers to excluded cases without listing them; no exclusions beyond the displayed inequality have been added.
\subsection{Short English statement}
The question asks whether every allowed surface type with Euler characteristic less than \(-1\) occurs as the fiber of a hyperbolic mapping torus whose trace field embeds into \(\R\). The explicit formulation uses stated genus, puncture, orientability, hyperbolicity, and trace-field conventions that the supplied record does not fully specify. It is a statement-only draft requiring mathematical review.
\subsection{Sources}
\begin{itemize}
\item Local UnsolvedMath registry, supplied record with permanent ID \(30001542\), title \emph{Real Places of Mapping Torus Trace Fields}, exact source problem code \texttt{catalogue entry}, in the Oberwolfach Reports Open Problems collection. The proposed year is recorded as 2010; no individual proposer is supplied. The registry labels the statement \texttt{corrected\_verified}, explaining that the parameters and inequality were restored to mathematical notation without changing the question. This is supplied verification metadata, not independent verification here.
\item Supplied report citation: \emph{Low-Dimensional Topology and Number Theory} (2011), Oberwolfach Reports, pp.\ 2101--2163, DOI \texttt{10.4171/owr/2010/35}. Original OWR report: \url{https://doi.org/10.4171/owr/2010/35}. The supplied background describes these reports as records of mathematical workshops and their open-problem discussions.
\item Supplied literature reference, labeled ``hyperbolic trace-field study'': \url{https://arxiv.org/abs/1609.08719}. No authors, article title, or specific supporting theorem are supplied in the record.
\item Supplied literature assessment dated exactly \texttt{2026-08-21}: the registry labels the problem partially solved and reports that the conjecture is proved for several infinite families, but not in full generality. The listed literature sources are the original OWR report and the hyperbolic trace-field study linked above. This assessment is supplied dated curation derived from the OWR material, not independently verified here; its time-sensitive status and cited claims require independent checking.
\end{itemize}
\end{document}
